\documentclass[11pt, a4paper]{article}

\usepackage[top = 1 in, bottom = 1 in, left = 0.5 in, right = 0.5 in]{geometry}

\usepackage{amsmath, amssymb, amsfonts}
\allowdisplaybreaks[4]

\usepackage{enumerate}
\usepackage{multirow}
\usepackage{hhline}
\usepackage{array}
\usepackage{longtable}
\usepackage{graphicx}
\usepackage{tabularray}
\usepackage{undertilde}
\usepackage{dingbat}
\usepackage{fontawesome5}
\usepackage[colorlinks=true, linkcolor=blue, urlcolor=red]{hyperref}
\usepackage{tasks}
\usepackage{bbding}
\usepackage{twemojis}
% how to use bull's eye ----- \scalebox{2.0}{\twemoji{bullseye}}
\usepackage{customdice}
% how to put dice face ------ \dice{2}
\usepackage{fontspec}
\usepackage{fancyhdr}
\usepackage{lastpage}
\usepackage{skak}
\usepackage{figchild}

\usepackage{soul}
\usepackage{xcolor}


\pagestyle{fancy}
\fancyhf{}
\fancyfoot[C]{Page \thepage\ of \pageref{LastPage}}
\renewcommand{\headrulewidth}{0pt}



\title{MSMS 407 : Longitudinal Data Analysis}
\author{{\fontsize{25}{20} \benyorithafont Ananda Biswas}}
\date{}

\newfontface\gloriahallelujahfont{GLORIAHALLELUJAH-REGULAR.TTF}

\newfontface\benyorithafont{Benyoritha-G334D.otf}

\newcommand{\bs}[1]{\boldsymbol{#1}}

\newcommand{\yhl}[1]{%
    {\sethlcolor{yellow}\hl{#1}}%
}

\newcommand{\rhl}[1]{%
    {\sethlcolor{red}\hl{#1}}%
}

\begin{document}

\maketitle

\section*{\faArrowAltCircleRight[regular] \textcolor{blue}{Restricted Maximum Likelihood Estimation}}

The method of REML estimation, was introduced by \textbf{Patterson} and
\textbf{Thompson} $(1971)$ as a way of \yhl{estimating variance components} in a GLM. The \rhl{objection} to the standard maximum likelihood procedure is that it produces \yhl{biased estimators of the covariance parameters}. This is well known in the case of the GLM with independent errors,

\begin{equation}\label{Y_independent_errors}
\boldsymbol{Y} \sim MVN(X\boldsymbol{\beta}, \sigma^2 I).
\end{equation}

In this case, the maximum likelihood estimator of $\sigma^2$ is

\begin{align*}
\widehat{\sigma^2} = \dfrac{\text{RSS}}{mn},
\end{align*}

where RSS denotes the residual sum of squares, whereas the usual unbiased estimator is

\begin{align*}
\widetilde{\sigma^2} = \dfrac{\text{RSS}}{mn - p},
\end{align*}

where $p$ is the number of elements of $\boldsymbol{\beta}$. In fact, $\widetilde{\sigma^2}$ is the REML estimator for $\sigma^2$ in the model (\ref{Y_independent_errors}). \\[0.5em]

In the case of the GLM with dependent errors,

\begin{equation}\label{Y_dependent_errors}
\boldsymbol{Y} \sim MVN(X\boldsymbol{\beta}, \sigma^2 V),
\end{equation}

\yhl{the REML estimator is defined as a maximum likelihood estimator based on a linearly transformed set of data $\boldsymbol{Y}^* = A\boldsymbol{Y}$ such that the distribution of $\boldsymbol{Y}^*$ does not depend on $\boldsymbol{\beta}$}. One way to achieve this is by taking $A$ to be the matrix which converts $\boldsymbol{Y}$ to OLS residuals (see Appendix \ref{app:Y_to_residuals}),

\begin{equation}\label{A}
A = I - X(X'X)^{-1}X'.
\end{equation}

Then, $\bs{Y}^{*}$ has a \yhl{singular} multivariate Gaussian distribution with \yhl{mean zero}, whatever the value of $\bs{\beta}$ (see Appendix \ref{app:Y_star_singular}). \\[0.5em]

To obtain a \yhl{non-singular distribution}, we could use only $nm - p$ rows of the matrix, $A$, defined at (\ref{A}) (see Appendix \ref{app:error_contrast}). It turns out that the resulting estimators for $\sigma^2$ and $V$ do not depend on which rows we use, nor indeed on the particular choice of $A$: any full-rank matrix with the property that $\mathbb{E}(\bs{Y}^{*}) = 0$ for all $\bs{\beta}$ will give the same answer. \\[0.5em]

Furthermore, from an operational viewpoint we do not need to make the transformation from $\bs{Y}$ to $\bs{Y}^{*}$ explicit, for the final estimators do not depend on $A$. The algebraic details are as follows. \\[0.5em]

For the theoretical development it is convenient to re-absorb $\sigma^2$ into $V$, and write our model for the response vector $\bs{Y}$ as

\begin{align}\label{distribution_of_Y_general}
\bs{Y} \sim MVN(X\beta, H),
\end{align}

where $H \equiv H(\bs{\alpha})$, the variance matrix of $\bs{Y}$, is characterized by a vector of parameters $\bs{\alpha}$, for example, in Uniform correlation model, $\bs{\alpha} = (\sigma^2, \rho)$. \\[0.5em]

Our \yhl{plan} is to transform $\bs{Y}$ to $\bs{Z}$ whose distribution does not depend of $\bs{\beta}$ and use the likelihood of $\bs{Z}$ to estimate the covariance parameters $\bs{\alpha}$. \\[0.5em]

Let $A$ be the matrix defined in (\ref{A}), and $B$ the $nm \times (nm - p)$ matrix defined by the requirements that $BB' = A$ and $B'B = I$, where $I$ denotes the $(nm - p) \times (nm - p)$ identity matrix. Finally, let
\[
\bs{Z}_{(nm - p) \times 1} = B'\bs{Y}.
\]

Note that,

\begin{align*}
\mathbb{E}(\bs{Z}) &= B'\mathbb{E}(\bs{Y}) \\[0.5em]
&= B'X\bs{\beta} \\[0.5em]
&= \underbrace{B'B}_{I}B'X\bs{\beta}, \; \text{since } B'B = I \\[0.5em]
&= B'\underbrace{BB'}_{A}X\bs{\beta} \\[0.5em]
&= B'AX\bs{\beta} \\[0.5em]
&= B' \left[ I - X(X'X)^{-1}X' \right]X\bs{\beta} \\[0.5em]
&= B' \left[ X - X(X'X)^{-1}X'X \right] \bs{\beta} \\[0.5em]
&= B' \left[ X - X \right] \bs{\beta} \\[0.5em]
&= B' \bs{O} B \\[0.5em]
&= \bs{0}.
\end{align*}

$\bs{Y} \sim MVN(X\beta, H)$, then $\bs{Z} = B'\bs{Y} \sim MVN(\bs{0}, B'HB)$ and the REML likelihood is

\begin{align}\label{Likelihood_based_on_Z}
L^{*}(\bs{\alpha})
=
\dfrac{1}{(2\pi)^{(nm-p)/2}\lvert B'HB \rvert^{1/2}}
\exp\!\left\{
-\dfrac{1}{2}
\bs{z}^{T}(B'HB)^{-1}\bs{z}
\right\}.
\end{align}

It can be clearly seen that the resulting estimator $\widehat{\bs{\alpha}}$ will depend on the arbitrary matrix $B$, hence different choices of $B$ will give different estimators. Also, this representation hides the connection with GLS residuals. So proceeding with the likelihood (\ref{Likelihood_based_on_Z}) is not the greatest of ideas. We take a different route. \\[0.5em]

First note that, for fixed $\bs{\alpha}$ the maximum likelihood estimator for $\bs{\beta}$ is the generalized least-squares estimator,
\[
\widehat{\bs{\beta}}
=
(X'H^{-1}X)^{-1}X'H^{-1}\bs{Y}
=
G\bs{Y}, \; \text{say},
\]

and

\begin{align}\label{distribution_of_beta_hat}
\widehat{\bs{\beta}} \sim MVN(\bs{\beta}, (X'HX)^{-1}).
\end{align}

Secondly recall, $$\bs{Z} = B'\bs{Y}.$$

Consider the transformation

\begin{align*}
\bs{Y}
\mapsto
\begin{pmatrix}
\bs{Z} \\
\widehat{\bs{\beta}}
\end{pmatrix}
=
\begin{pmatrix}
B' \\
G
\end{pmatrix}
\bs{Y}
= T\bs{Y}, \; \text{say }.
\end{align*}

We write,

\begin{align*}
f_{\bs{Z},\widehat{\bs{\beta}}}(\bs{z},\widehat{\bs{\beta}})
=
f_{\bs{Y}}(\bs{y}) \cdot J
\end{align*}

where $J$ is the Jacobian of the transformation given by

\begin{align*}
J
=
\dfrac{1}{
\left\lvert
\dfrac{\partial}{\partial \bs{Y}}\, T\bs{Y}
\right\rvert
}
= \dfrac{1}{\left\lvert T \right\rvert}
= \left\lvert T \right\rvert^{-1}.
\end{align*}

$\bs{Z}$ and $\widehat{\bs{\beta}}$ are \yhl{independently distributed} (see Appendix \ref{app:Z_and_beta_independent}) and we can write

\begin{align*}
f_{\bs{Z}}(\bs{z}) \cdot f_{\widehat{\bs{\beta}}}(\widehat{\bs{\beta}}) 
&= f_{\bs{Y}}(\bs{y}) \cdot J \\[1em]
\Rightarrow f_{\bs{Z}}(\bs{z})
&=
\dfrac{
f_{\bs{Y}}(\bs{y})
}{
f_{\widehat{\bs{\beta}}}(\widehat{\bs{\beta}})}
\cdot 
\left\lvert T \right\rvert^{-1}
.
\end{align*}

The Jacobian is free from the parameters $\bs{\beta}$ and $\bs{\alpha}$ (see Appendix \ref{app:determinant_of_T}) and thus can be ignored while making inferences. \\[0.5em]

Finally, we can write

\begin{align}\label{PDF_Z_form}
f_{\bs{Z}}(\bs{z})
\propto
\dfrac{
f_{\bs{Y}}(\bs{y})
}{
f_{\widehat{\bs{\beta}}}(\widehat{\bs{\beta}})
}.
\end{align}

From (\ref{distribution_of_Y_general}),

\begin{align}\label{PDF_Y_general}
f_{\bs{y}}(\bs{y})
=
\dfrac{
1
}{
(2\pi)^{nm/2}
\lvert H \rvert^{1/2}
}
\exp\!\left\{
-\dfrac{1}{2}
\left(
\bs{y} - X\bs{\beta}
\right)'
H^{-1}
\left(
\bs{y} - X\bs{\beta}
\right)
\right\}.
\end{align}

From (\ref{distribution_of_beta_hat}),

\begin{align}\label{PDF_beta_hat}
f_{\widehat{\bs{\beta}}}(\widehat{\bs{\beta}})
=
\dfrac{
1
}{
(2\pi)^{p/2}
\left\lvert
X'H^{-1}X
\right\rvert^{-1/2}
}
\exp\!\left\{
-\dfrac{1}{2}
\left(
\widehat{\bs{\beta}} - \bs{\beta}
\right)'
\left(
X'H^{-1}X
\right)
\left(
\widehat{\bs{\beta}} - \bs{\beta}
\right)
\right\}.
\end{align}


\newpage

We state a result,

\begin{align} \label{result}
\left(
\bs{y} - X\bs{\beta}
\right)'
H^{-1}
\left(
\bs{y} - X\bs{\beta}
\right)
=
\left(
\bs{y} - X\widehat{\bs{\beta}}
\right)'
H^{-1}
\left(
\bs{y} - X\widehat{\bs{\beta}}
\right)
+
\left(
\widehat{\bs{\beta}} - \bs{\beta}
\right)'
\left(
X'H^{-1}X
\right)
\left(
\widehat{\bs{\beta}} - \bs{\beta}
\right).
\end{align}

Using (\ref{PDF_Z_form}), (\ref{PDF_Y_general}), (\ref{PDF_beta_hat}) along with (\ref{result}), we get the PDF of $\bs{Z} = B'\bs{Y}$ is proportional to 

\begin{align}\label{REML_likelihood}
\dfrac{f_{\bs{y}}(\bs{y})}{f_{\widehat{\bs{\beta}}}(\widehat{\bs{\beta}})} =
(2\pi)^{-\frac{1}{2}(nm-p)}
\lvert H \rvert^{-1/2}
\left\lvert
X'H^{-1}X
\right\rvert^{-1/2}
\times
\exp\!\left\{
-\dfrac{1}{2}
\left(
\bs{y} - X\widehat{\bs{\beta}}
\right)'
H^{-1}
\left(
\bs{y} - X\widehat{\bs{\beta}}
\right)
\right\}
\end{align}

Note that the right-hand side of (\ref{REML_likelihood}) is independent of $A$, and the same result would therefore hold for any $\bs{Z}$ such that $\mathbb{E}(\bs{Z}) = \bs{0}$ and $\operatorname{Cov}(\bs{Z}, \widehat{\bs{\beta}}) = \bs{0}$. \\[0.5em]

The practical implication of (\ref{REML_likelihood}) is that the REML estimator, $\widetilde{\bs{\alpha}}$, maximizes the log-likelihood

\begin{align}\label{REML_log_likelihood}
\ell^{*}(\bs{\alpha})
=
-\dfrac{1}{2}
\log \lvert H \rvert
-
\dfrac{1}{2}
\log
\left\lvert
X'H^{-1}X
\right\rvert
-
\dfrac{1}{2}
\left(
\bs{y} - X\widehat{\bs{\beta}}
\right)'
H^{-1}
\left(
\bs{y} - X\widehat{\bs{\beta}}
\right),
\end{align}

whereas the maximum likelihood estimator $\widehat{\bs{\alpha}}$ maximizes

\[
\ell(\bs{\alpha})
=
-\dfrac{1}{2}
\log \lvert H \rvert
-
\dfrac{1}{2}
\left(
\bs{y} - X\widehat{\bs{\beta}}
\right)'
H^{-1}
\left(
\bs{y} - X\widehat{\bs{\beta}}
\right).
\]

It follows from this last result that the algorithm to implement REML
estimation for the model (\ref{Y_dependent_errors}) incorporates only a simple modification to the maximum likelihood algorithm derived earlier. \\[0.5em]

\fcAntA[scale = 0.4, draw = red, line width = 1] Following (\ref{Y_dependent_errors}), recall that we consider $m$ units with $n$ measurements per unit and that $\sigma^2V$ is a block-diagonal matrix with non-zero $n \times n$ blocks $\sigma^2V_0$ representing the common covariance matrix of the measurements on any one unit. \\[0.5em]

The original complete log-likelihood is

\begin{align}\label{original_log_likelihood}
\ell(\boldsymbol{\beta}, \sigma^2, V_0) &= -\dfrac{1}{2} \; mn \ln(\sigma^2) - \dfrac{1}{2} \; m \ln(\left| V_0 \right|) - \dfrac{1}{2\sigma^2} \; (\boldsymbol{y} - X\boldsymbol{\beta})' V^{-1} (\boldsymbol{y} - X\boldsymbol{\beta}).
\end{align}

In (\ref{original_log_likelihood}), for given $V_0$, $V$ will be completely known and we write

\begin{align}\label{beta_hat}
\widehat{\boldsymbol{\beta}}(V_0) = (X'V^{-1}X)^{-1}X'V^{-1}\boldsymbol{y}.
\end{align}

and

\begin{align}\label{RSS}
\text{RSS}(V_0) &= \left\{\boldsymbol{y} - X\widehat{\boldsymbol{\beta}}(V_0)\right\}' V^{-1} \left\{\boldsymbol{y} - X\widehat{\boldsymbol{\beta}}(V_0)\right\}.
\end{align}

Put $H(\bs{\alpha}) = \sigma^2 V$ in (\ref{distribution_of_Y_general}) with $\bs{\alpha} = (\sigma^2, V_0)$. On using (\ref{beta_hat}) and $H(\bs{\alpha}) = \sigma^2 V$, general (\ref{REML_log_likelihood}) becomes special case

\begin{align}\label{REML_likelihood_special_case}
\ell^{*}(\sigma^2, V_0)
&=
-\dfrac{1}{2}
\log \lvert \sigma^2 V \rvert
-
\dfrac{1}{2}
\log
\left\lvert
X'(\sigma^2 V)^{-1}X
\right\rvert
-
\dfrac{1}{2}
\left(
\bs{y} - X\widehat{\bs{\beta}}
\right)'
(\sigma^2 V)^{-1}
\left(
\bs{y} - X\widehat{\bs{\beta}}
\right) \nonumber \\[1.5em]
&=
-\dfrac{1}{2}
\log \left\{ (\sigma^2)^{mn} \lvert V \rvert \right\}
-
\dfrac{1}{2}
\log \left\{
\left( \dfrac{1}{\sigma^2} \right)^p
\left\lvert
X'V^{-1}X
\right\rvert
\right\}
-
\dfrac{1}{2\sigma^2}
\left(
\bs{y} - X\widehat{\bs{\beta}}
\right)'
V^{-1}
\left(
\bs{y} - X\widehat{\bs{\beta}}
\right) \nonumber \\[0.5em]
&=
-\dfrac{1}{2}(mn-p)\log \sigma^2
-
\dfrac{1}{2} m \cdot \log \lvert V_0 \rvert
-
\dfrac{1}{2}
\log
\left\lvert
X^{T}V^{-1}X
\right\rvert
-
\dfrac{1}{2\sigma^2}
\left(
\bs{y}-X\widehat{\bs{\beta}}
\right)^{T}
V^{-1}
\left(
\bs{y}-X\widehat{\bs{\beta}}
\right)
\end{align}



Then maximizing (\ref{REML_likelihood_special_case}) w.r.t. $\sigma^2$ gives the REML estimator for $\sigma^2$ as

\begin{align}\label{sigma_sq_hat_REML}
\widetilde{\sigma}^{2}(V_0)
=
\dfrac{
\text{RSS}(V_0)
}{
nm-p
},
\end{align}

where $p$ is the number of elements of $\bs{\beta}$. The REML estimator for $V_0$ maximizes the reduced REML log-likelihood

\begin{align}\label{Reduced_REML_log_likelihood}
\ell^{*}(V_0)
=
-\dfrac{1}{2}
m
\left\{
n\log \text{RSS}(V_0)
+
\log\!\left(
\lvert V_0 \rvert
\right)
\right\}
-
\dfrac{1}{2}
\log\!\left(
\left\lvert
X'V^{-1}X
\right\rvert
\right).
\end{align}


Finally, maximization of (\ref{Reduced_REML_log_likelihood}) w.r.t. distinct elements of $V_0$ yields $\widetilde{V}_0$. Substitution of the resulting estimator $\widetilde{V}_0$ into (\ref{beta_hat}) and (\ref{sigma_sq_hat_REML}) gives the REML estimators

\[
\widetilde{\bs{\beta}}
=
\widehat{\bs{\beta}}(\widetilde{V}_0)
\quad \& \quad
\widetilde{\sigma}^{2}
=
\widetilde{\sigma}^{2}(\widetilde{V}_0).
\]

\section*{Appendix}

\appendix

\section{$\boldsymbol{Y}$ to OLS residuals}
\label{app:Y_to_residuals}

In OLS, 
$$\hat{\bs{\beta}} = (X'X)^{-1}X'\bs{Y},$$
\vspace{0.2cm}
$$\widehat{\bs{Y}} = X(X'X)^{-1}X'\bs{Y}.$$

Then,

\begin{align*}
\bs{e} &= \bs{Y} - \widehat{\bs{Y}} \\[0.5em]
&= \bs{Y} - X(X'X)^{-1}X'\bs{Y} \\[0.5em]
&= \left[ I - X(X'X)^{-1}X' \right] \bs{Y}.
\end{align*}

\section{$\bs{Y}^*$ has a singular distribution with mean $\bs{0}$}
\label{app:Y_star_singular}

$\bs{Y}^* = A\bs{Y}$, $\bs{Y}$ has a multivariate normal distribution, consequently $\bs{Y}^*$ also has a multivariate normal distribution.

$$E(\boldsymbol{Y}^*) = A \cdot E(\boldsymbol{Y}) = (I - X(X'X)^{-1}X') \cdot X\boldsymbol{\beta} = X\boldsymbol{\beta} - X(X'X)^{-1}X'X\boldsymbol{\beta} = X\boldsymbol{\beta} - X\boldsymbol{\beta} = \boldsymbol{0}$$

and

$$\text{disp}(\boldsymbol{Y}^*) = \sigma^2 \cdot AVA'.$$

Then $\left\lvert \text{disp}(\bs{Y}^*) \right\rvert = \sigma^2 \cdot \left\lvert V \right\rvert \left\lvert A \right\rvert^2$. \\[0.5em]

Also, $AX = \left[ I - X(X'X)^{-1}X' \right]X = \left[ X - X(X'X)^{-1}X'X \right] = \left[ X - X \right] = \bs{O}.$ 

$$\therefore \; A\utilde{x} = \utilde{0} \Rightarrow \utilde{x} \in \mathcal{N}(A) \Rightarrow \text{Nullity}(A) \neq 0 \Rightarrow \text{rank}(A) < mn \Rightarrow \left\lvert A \right\rvert = 0.$$

$$\therefore \; \left\lvert \text{disp}(\bs{Y}^*) \right\rvert = 0 \Rightarrow \text{disp}(\bs{Y}^*) \text{ has a singular distribution.}$$



\section{Error Contrast}
\label{app:error_contrast}

A linear combination $\bs{a}'\bs{Y}$ of the response variable $\bs{Y}$ is called an \textbf{error contrast} if $\mathbb{E}(\bs{a}'\bs{Y}) = 0$. For a linear model $\bs{Y} = X\bs{\beta} + \bs{\epsilon}$, the condition reduces to $\bs{a}'X = \bs{0}$. \\[0.5em]

It readily follows that each row of $A$ as in (\ref{A}) is an error contrast as $AX = \bs{O}$. Recall the fact that, given $n$ error contrasts, at max $n-p$ can be linearly independent, $p$ being the rank of the related design matrix $X$. $\bs{Y^*} = A\bs{Y}$ is then a vector of error contrasts as each row of $A$ produces an error contrast. Thus to have a nonsingular distribution of $\bs{Y^*}$ we must take $mn-p$ rows of $A$, of course the $mn-p$ rows must be linearly independent.

\section{$\bs{Z}$ and $\bs{\widehat{\beta}}$ are independent}
\label{app:Z_and_beta_independent}

\begin{align*}
	\operatorname{Cov}(\bs{Z},\widehat{\bs{\beta}})
	&= \mathbb{E}\left\{\bs{Z}\left(\widehat{\bs{\beta}}-\bs{\beta}\right)'\right\} \\
	&= \mathbb{E}\left\{B'\bs{Y}\left(\bs{Y}'G'-\bs{\beta}'\right)\right\} \\
	&= B'\mathbb{E}\left(\bs{Y}\bs{Y}'\right)G'
	- B'\mathbb{E}\left(\bs{Y}\right)\bs{\beta}' \\
	&= B'\left\{\operatorname{Var}(\bs{Y})
	+ \mathbb{E}(\bs{Y})\mathbb{E}(\bs{Y})'\right\}G'
	- B'\mathbb{E}(\bs{Y})\bs{\beta}' .
\end{align*}

Now, substituting $\mathbb{E}(\bs{Y}) = X\bs{\beta}$ into this last expression gives

\begin{align*}
	\operatorname{Cov}(\bs{Z},\widehat{\bs{\beta}})
	&= B'\left(H + X\bs{\beta}\bs{\beta}'X'\right)G'
	- B'X\bs{\beta}\bs{\beta}' .
\end{align*}

Also,

\begin{align*}
	X'G'
	&= X'H^{-1}X\left(X'H^{-1}X\right)^{-1} \\
	&= I,
\end{align*}

and

\begin{align*}
	B'HG'
	&= B'HH^{-1}X\left(X'H^{-1}X\right)^{-1} \\
	&= B'X\left(X'H^{-1}X\right)^{-1} \\
	&= 0,
\end{align*}

because $B'X = B'AX = \bs{O}$ as in the proof that
$\mathbb{E}(\bs{Z}) = \bs{0}$. It follows that

\begin{align*}
	\operatorname{Cov}(\bs{Z},\widehat{\bs{\beta}})
	&= B'X\bs{\beta}\bs{\beta}'
	- B'X\bs{\beta}\bs{\beta}' \\
	&= \bs{O}.
\end{align*}

Finally, in the multivariate Gaussian setting, zero covariance is equivalent to independence.

\section{$\mid T \mid$}
\label{app:determinant_of_T}

\begin{align*}
	\mid T \mid &= \mid TT' \mid^{1/2} \\[0.2em]
	&= \left[\begin{pmatrix}
		B' \\
		G
	\end{pmatrix}
	\begin{pmatrix}
		B \; G'
	\end{pmatrix}\right]^{1/2} \\[0.2em]
	&= \begin{vmatrix}
		B'B & B'G' \\
		GB & GG'
	\end{vmatrix}^{1/2} \\[0.2em]
	&= \begin{vmatrix}
		I & B'G' \\
		GB & GG'
	\end{vmatrix}^{1/2} \; \text{as } B'B = I \\[0.2em]
	&= \mid I \mid^{1/2} \cdot \mid GG' - GBI^{-1}B'G' \mid^{1/2} \\[0.2em]
	&= \mid GG' - GBB'G' \mid^{1/2} \\[0.2em]
	&= \mid GG' - GAG' \mid^{1/2} \; \text{as } BB' = A \\[0.2em]
	&= \mid GG' - G\left[I - X(X'X)^{-1}X'\right]G' \mid^{1/2} \\[0.2em]
	&= \mid GG' - \left[GG' - GX(X'X)^{-1}X'G'\right] \mid^{1/2} \\[0.2em]
	&= \mid GG' - GG' + GX(X'X)^{-1}X'G' \mid^{1/2} \\[0.2em]
	&= \mid GX(X'X)^{-1}X'G' \mid^{1/2} \\[0.2em]
	&= \mid (X'X)^{-1} \mid^{1/2} \; \text{as } GX = I \\[0.2em]
	&= \mid X'X \mid^{-1/2}.
\end{align*}

\end{document}